\chapter{Introduzione all'informatica}
\label{chap:introduzione-informatica}

L'informatica studia il trattamento automatico dell'informazione e i
processi di calcolo che lo rendono possibile. Il suo oggetto non è
soltanto il calcolatore come dispositivo fisico, ma anche la formulazione
dei problemi, la progettazione dei procedimenti risolutivi e i limiti
della loro esecuzione. Per comprendere come un programma trasformi i
dati in risultati occorre distinguere il problema da risolvere, il
metodo scelto e la macchina che lo esegue.\footnote{Riferimento didattico:
S. Canazza, \emph{Presentazione del corso}, file
\nolinkurl{2026_FI_00_canazza.pdf}.}

\section{Informazione, calcolo e programmazione}
\label{sec:fi-informatica-calcolo}

Un \emph{dato} è una rappresentazione utilizzabile per conservare,
trasmettere o elaborare informazione. Il suo significato dipende dalle
convenzioni adottate: una sequenza di simboli può descrivere, per esempio,
un importo oppure un testo. L'elaborazione opera sulle rappresentazioni,
mentre l'interpretazione dei risultati deve rimanere coerente con il
problema originario. Le convenzioni per rappresentare numeri e caratteri
sono sviluppate nel \cref{chap:rappresentazione-informazioni}.

Il significato di un dato non coincide, tuttavia, con la quantità di
informazione che esso può trasmettere. Nel modello introdotto da Claude
Shannon, la quantità di informazione associata a un evento dipende dalla
sua probabilità, indipendentemente dal significato del messaggio.
Un evento poco probabile fornisce, quando si verifica, una quantità
di informazione maggiore rispetto a uno più
probabile. Per un esito di probabilità \(0<p\le1\), il contenuto
informativo è
\[
  I(p)=-\log_2 p.
\]
Il logaritmo in base due esprime questa quantità in \emph{bit}: un esito
fra due alternative equiprobabili fornisce un bit, mentre un esito certo
fornisce zero bit. Per esempio, conoscere l'esito di un esperimento con
quattro risultati equiprobabili fornisce \(-\log_2(1/4)=2\) bit.
Per esiti con probabilità \(p_1,\ldots,p_r\), la cui somma è uno,
l'\emph{entropia di Shannon} misura il contenuto informativo medio:
\[
  H=-\sum_{i=1}^{r}p_i\log_2 p_i,
\]
con contributo nullo per gli esiti di probabilità zero. Il bit come unità
di informazione non coincide concettualmente con una \emph{cifra binaria}:
quest'ultima è un simbolo \(0\) o \(1\). Una cifra binaria aleatoria
ha entropia di un bit solo se i due valori sono equiprobabili; se il
valore è già noto, la sua entropia è nulla.\footnote{C. E. Shannon,
\emph{A Mathematical Theory of Communication}, \emph{Bell System
Technical Journal}, 1948, parte I: misura dell'informazione ed entropia
per sorgenti discrete.}

L'elaborazione automatica dei dati è affidata a un \emph{computer},
o \emph{calcolatore}, ossia un dispositivo fisico programmabile che
memorizza dati, esegue istruzioni e scambia informazioni con l'esterno.
Una \emph{computazione} è invece un processo di trasformazione dei dati
secondo regole specificate. Il dispositivo e il processo sono concetti
distinti: lo stesso calcolatore può eseguire computazioni diverse,
e lo stesso procedimento può essere realizzato su macchine differenti.
Le componenti fisiche costituiscono l'\emph{hardware}, mentre i
programmi costituiscono il \emph{software}.

Per stabilire quale elaborazione debba essere eseguita, occorre
anzitutto definire il \emph{problema di calcolo}, specificando i dati
ammissibili in ingresso e il risultato richiesto per ciascuno di essi.
Una particolare scelta dei dati costituisce un'\emph{istanza} del
problema. La specifica chiarisce che cosa ottenere, ma non come
ottenerlo: il procedimento risolutivo è un \emph{algoritmo}, mentre
un \emph{programma} ne esprime una realizzazione in un linguaggio di
programmazione. Uno stesso algoritmo può avere più implementazioni,
differenti per linguaggio o organizzazione interna, senza modificare
il problema risolto.

La distinzione fra problema, algoritmo e programma permette di precisare
il ruolo della \emph{programmazione}, che comprende la progettazione,
l'implementazione e la verifica dei programmi rispetto alle specifiche.
Essa richiede di individuare rappresentazioni dei dati e procedimenti
risolutivi appropriati, traducendoli in istruzioni conformi al linguaggio
utilizzato. Il calcolatore esegue tali istruzioni secondo le regole
del linguaggio e dell'ambiente di esecuzione, senza correggerle
autonomamente sulla base delle intenzioni del programmatore.

A livello della macchina, anche i procedimenti più articolati si
riconducono a operazioni elementari, come trasferimenti di dati,
calcoli, confronti e scelte della prossima istruzione. La loro
combinazione mediante sequenze, selezioni fra percorsi alternativi
e ripetizioni permette di realizzare elaborazioni complesse a partire
da un insieme limitato di operazioni fondamentali.

\section{Algoritmi e risoluzione dei problemi}
\label{sec:fi-algoritmi}

La descrizione di un metodo deve precisare tanto le operazioni quanto
il loro ordine e le condizioni che determinano il proseguimento del
calcolo. Un'indicazione come «ripetere finché il risultato è adeguato»
non è sufficiente se non viene stabilito un criterio verificabile di
adeguatezza. Analogamente, chiedere di individuare il paesaggio «più
rilassante» richiede prima di definire un criterio preciso di confronto.

\begin{definition}[Algoritmo per un problema di calcolo]
\label{def:fi-algoritmo}
Un algoritmo è una descrizione finita e non ambigua di un procedimento,
composto da passi effettivamente eseguibili, che termina dopo un numero
finito di passi per ogni ingresso del dominio previsto. Per risolvere
correttamente un problema, il procedimento deve inoltre produrre il
risultato prescritto per ciascun ingresso ammissibile.
\end{definition}

L'\emph{eseguibilità} richiede che ogni passo possa essere svolto mediante
operazioni definite, senza richiedere informazioni non disponibili o
procedure di cui non sia specificato il funzionamento.
L'\emph{assenza di ambiguità} riguarda sia il significato delle operazioni
sia la scelta del passo successivo. La \emph{terminazione} riguarda
ogni ingresso ammissibile: una descrizione di poche istruzioni può
imporre una ripetizione senza fine, quindi la finitezza del testo non
assicura la finitezza dell'esecuzione.

La terminazione non impone lo stesso numero di passi per tutti gli
ingressi, né garantisce che il risultato sia corretto. Correttezza e
terminazione devono essere giustificate rispetto alle ipotesi del
problema. Un programma che implementa un algoritmo deve inoltre
rispettarne tali ipotesi; errori nella traduzione o limiti della
rappresentazione dei dati possono alterare il risultato.

Un semplice problema numerico permette di esaminare insieme
terminazione e correttezza di un procedimento iterativo.

\noindent\begin{minipage}{\linewidth}
\begin{example}[Raddoppio di un capitale]
\label{ex:fi-interesse-composto}
Si vuole determinare il primo anno in cui un capitale di \(20\,000\) euro
raggiunge almeno \(40\,000\) euro, con interesse annuo costante del \(5\%\)
capitalizzato a fine anno. Si escludono versamenti, prelievi, imposte e
arrotondamenti intermedi. Il procedimento è il seguente:
\begin{enumerate}
  \item inizializzare il numero di anni \(n=0\) e il saldo \(S=20\,000\);
  \item finché \(S<40\,000\), sostituire \(S\) con \(1{,}05S\) e aumentare
        \(n\) di uno;
  \item restituire \(n\).
\end{enumerate}
Dopo \(n\) ripetizioni il saldo è \(S_n=20\,000(1{,}05)^n\): ogni passo
aggiunge al saldo il \(5\%\) del suo valore corrente. Finché la soglia
non viene raggiunta, il saldo è almeno \(20\,000\) euro e aumenta di
almeno \(1\,000\) euro all'anno. La soglia viene quindi raggiunta in
non più di venti ripetizioni, garantendo la terminazione.

Poiché la condizione viene controllata prima di ogni aggiornamento,
il valore restituito corrisponde al primo anno in cui il capitale
raggiunge la soglia. In particolare,
\[
  S_{14}\approx39\,598{,}63<40\,000,
  \qquad S_{15}\approx41\,578{,}56>40\,000,
\]
e il risultato è \(15\) anni. Le approssimazioni mostrate non sostituiscono
l'ipotesi di calcolo esatto del modello.
\end{example}
\end{minipage}\par

L'\cref{ex:fi-interesse-composto} mostra come sia possibile descrivere
un procedimento e giustificarne terminazione e correttezza senza
ricorrere a un particolare linguaggio di programmazione. La sua
implementazione richiede invece di tradurre inizializzazione,
confronto, aggiornamento e ripetizione nelle istruzioni previste
dal linguaggio scelto.

Il controllo di un algoritmo può anche essere rappresentato mediante
un diagramma di flusso. Quello seguente illustra un semplice conteggio.

\noindent\begin{minipage}[c]{0.40\linewidth}
  \centering
  \begin{tikzpicture}[unipd/diagram,node distance=7mm]
  \node[unipd/flow start] (inizio) {Inizio};
  \node[unipd/flow process,rounded corners=0pt,below=of inizio]
    (inizializzazione) {\(n\gets0\)};
  \node[unipd/flow decision,below=of inizializzazione]
    (condizione) {\(n<3\)?};
  \node[unipd/flow process,rounded corners=0pt,below=10mm of condizione]
    (incremento) {\(n\gets n+1\)};
  \node[unipd/flow end,below=of incremento] (fine) {Fine};
  \draw[unipd/flow arrow] (inizio) -- (inizializzazione);
  \draw[unipd/flow arrow] (inizializzazione) -- (condizione);
  \draw[unipd/flow arrow] (condizione) --
    node[unipd/edge label,midway,right,xshift=1mm,yshift=1mm] {Sì} (incremento);
  \draw[unipd/flow arrow] (incremento.west) -- ++(-6mm,0)
    |- (condizione.west);
  \draw[unipd/flow arrow] (condizione.east) --
    node[unipd/edge label,midway,above] {No} ++(9mm,0)
    |- (fine.east);
\end{tikzpicture}

\end{minipage}\hfill
\begin{minipage}[c]{0.56\linewidth}
Un \emph{diagramma di flusso} rappresenta graficamente il controllo
di un algoritmo: ovali indicano inizio e fine, rettangoli le operazioni,
rombi le decisioni e frecce l'ordine dei passi. Le uscite da una
decisione corrispondono ai possibili esiti della condizione; una
freccia che ritorna a un passo precedente indica una ripetizione.
Il diagramma descrive quindi l'ordine delle operazioni, ma non
costituisce di per sé una prova della loro correttezza o terminazione.
\end{minipage}\par

Una volta precisato come descrivere un algoritmo e verificarne
le proprietà, occorre distinguere due questioni più generali:
l'esistenza di una soluzione algoritmica e le risorse necessarie
per eseguirla.

Un problema è \emph{calcolabile} se esiste un procedimento algoritmico
che lo risolve per tutti gli ingressi ammissibili. La calcolabilità
riguarda quindi l'esistenza del procedimento, indipendentemente
dal fatto che esso sia già noto. L'\emph{efficienza} riguarda invece
le risorse richieste dalla soluzione, soprattutto tempo e memoria,
al variare dei dati. La \emph{complessità computazionale} studia
queste risorse in funzione della dimensione dell'ingresso, secondo
un modello di calcolo specificato; non coincide necessariamente
con il tempo misurato su un particolare computer.

Anche un algoritmo che termina correttamente può richiedere risorse
troppo grandi per un impiego pratico. La possibilità teorica di
risolvere un problema e quella di farlo con le risorse disponibili
sono dunque questioni distinte.

\section{La macchina di Turing e i limiti del calcolo}
\label{sec:fi-macchina-turing}

La questione della calcolabilità, introdotta nella
\cref{sec:fi-algoritmi}, richiede un modello preciso del significato
di «procedimento eseguibile».
Nel contesto delle ricerche sui fondamenti della matematica promosse
da Hilbert, il problema della decisione chiedeva un metodo generale
per stabilire la validità delle formule della logica del primo ordine.
Nel 1936 Church e Turing ottennero, con approcci diversi, risultati
negativi su questa possibilità. Turing introdusse un modello di calcolo
astratto, oggi chiamato \emph{macchina di Turing}, che permette di
studiare il calcolo indipendentemente dalla tecnologia di una macchina
fisica.\footnote{Per il modello, il contesto storico, l'universalità e
il problema dell'arresto: L. De Mol,
\href{https://plato.stanford.edu/entries/turing-machine/}
{\emph{Turing Machines}},
\emph{Stanford Encyclopedia of Philosophy}, §§ 1--2 e 5.}

Nella versione deterministica qui considerata, la macchina dispone di
un nastro diviso in celle, ciascuna contenente un simbolo di un alfabeto
finito che comprende un simbolo di bianco. Il nastro è teoricamente
illimitato: non viene imposto un limite alla memoria utilizzabile.
L'ingresso occupa inizialmente un tratto finito, mentre le altre celle
sono bianche. Una testina legge e scrive una cella alla volta e può
spostarsi alla cella adiacente a destra o a sinistra.

Il controllo della macchina ha un insieme finito di \emph{stati}, fra
cui uno iniziale e uno di arresto. Le \emph{regole di transizione}
stabiliscono, in base allo stato corrente e al simbolo letto, quale
simbolo scrivere, dove spostare la testina e quale stato assumere.
Per ogni coppia di stato e simbolo per la quale è definita una regola,
l'azione è determinata senza ambiguità. Scrivere può anche significare
lasciare invariato il simbolo della cella.

\begin{figure}[htbp]
  \centering
  \begin{tikzpicture}[unipd/diagram]
  \foreach \i/\tapesymbol in {0/\mdlgwhtsquare,1/0,2/1,3/1,4/0,5/\mdlgwhtsquare,6/\mdlgwhtsquare} {
    \node[unipd/node,rectangle,rounded corners=0pt,
          minimum width=10mm,minimum height=10mm,inner sep=0pt]
      (cell\i) at (\i,0) {$\tapesymbol$};
  }
  \node at (-0.9,0) {$\cdots$};
  \node at (6.9,0) {$\cdots$};
  \node[unipd/node] (head) at (3,1.5) {Testina: stato $q$};
  \draw[unipd/directed edge,shorten <=2pt,shorten >=2pt]
    (head.south) -- (cell3.north);
  % Etichetta separata: il suo sfondo non deve coprire la freccia.
  \draw[unipd/edge,<->] (2,-0.9) -- (4,-0.9);
  \node[unipd/edge label] at (3,-1.4) {spostamento di una cella};
\end{tikzpicture}

  \caption{Un tratto del nastro e la posizione della testina.
    Il simbolo \(\mdlgwhtsquare\) indica il bianco; \(q\) è lo stato corrente.}
  \label{fig:fi-nastro-turing}
\end{figure}

La \cref{fig:fi-nastro-turing} mostra soltanto una parte del nastro,
non un limite alla sua estensione. Una \emph{configurazione} completa
comprende il contenuto del nastro, la posizione della testina e lo
stato del controllo. L'applicazione di una regola trasforma una
configurazione nella successiva; una computazione è la successione di
questi passi. Quando viene raggiunto lo stato di arresto, il risultato
si ricava dal nastro secondo la convenzione scelta. La macchina può
anche continuare indefinitamente. La finitezza dell'insieme degli stati
di controllo non impone un limite alla lunghezza del nastro; ogni
esecuzione che termina, tuttavia, visita soltanto un numero finito di celle.

Una \emph{macchina universale di Turing} riceve una descrizione codificata
di un'altra macchina e il suo ingresso, e ne simula l'esecuzione.
Le regole della macchina universale rimangono fisse: cambiando la
descrizione sul nastro cambia il procedimento simulato. Il programma
può quindi essere trattato come un dato. Universalità significa
capacità di simulare ogni macchina del modello, non capacità di risolvere
qualsiasi problema concepibile, né garanzia che la simulazione termini.

La \emph{tesi di Church--Turing} identifica la nozione intuitiva di
procedimento effettivo con la calcolabilità mediante macchina di Turing.
È una tesi sul significato del calcolo algoritmico, non un teorema
secondo cui ogni problema sarebbe risolvibile. Il modello fornisce
così un riferimento comune per distinguere problemi calcolabili e
problemi che non ammettono una soluzione algoritmica generale.

Fra questi rivestono particolare importanza i \emph{problemi di decisione},
che richiedono una risposta fra due alternative, convenzionalmente
«sì» e «no». Un problema di decisione è \emph{decidibile} se esiste
un algoritmo che termina per ogni ingresso ammissibile e fornisce
sempre la risposta corretta; in caso contrario è \emph{indecidibile}.
La decidibilità riguarda quindi l'esistenza di un procedimento
generale, non la possibilità di risolvere una singola istanza.

Un esempio fondamentale è il \emph{problema dell'arresto}: stabilire
se un programma, eseguito su un determinato ingresso, terminerà
oppure continuerà indefinitamente. Non basta osservare che una
particolare esecuzione termina; per decidere il problema occorrerebbe
un procedimento capace di fornire una risposta corretta in entrambi
i casi, per qualunque programma e ingresso.

\begin{proposition}[Indecidibilità del problema dell'arresto]
\label{prop:fi-problema-arresto}
Non esiste un algoritmo che, ricevuti la descrizione di una macchina di
Turing arbitraria e un suo ingresso, termini sempre e stabilisca
correttamente se quella macchina si arresterà su quell'ingresso.
\end{proposition}

\noindent\emph{Dimostrazione (per assurdo).}
Si supponga che esista un algoritmo, chiamato \emph{osservatore},
capace di stabilire se un qualsiasi programma terminerà su un dato
ingresso. Indichiamo con \(O(p,x)\) il risultato dell'osservatore,
dove \(p\) è la descrizione codificata di un programma e \(x\)
è il suo ingresso. Per ipotesi, l'osservatore termina sempre
e restituisce
\[
  O(p,x)=
  \begin{cases}
    1 & \text{se il programma descritto da \(p\) termina su \(x\)},\\
    0 & \text{se il programma descritto da \(p\) non termina su \(x\)}.
  \end{cases}
\]

L'ipotesi è più forte della semplice simulazione: l'osservatore
deve rispondere in tempo finito anche quando il programma esaminato
non si arresta.

Utilizzando l'osservatore, si costruisce un nuovo programma \(D\)
che riceve come ingresso la descrizione \(p\) di un programma.
Il programma \(D\) chiede all'osservatore se il programma descritto
da \(p\), eseguito sulla propria descrizione \(p\), terminerà.
Successivamente si comporta in modo opposto alla previsione:
\[
  D(p):
  \begin{cases}
    \text{esegue un ciclo senza fine}
      & \text{se } O(p,p)=1,\\
    \text{termina immediatamente}
      & \text{se } O(p,p)=0.
  \end{cases}
\]

La costruzione è possibile sotto l'ipotesi iniziale: \(D\) utilizza
l'osservatore come parte del proprio procedimento e sceglie quale
comportamento adottare in base alla risposta ricevuta.

Poiché anche \(D\) è un programma, possiede una descrizione codificata.
Indichiamola con \(d\) e consideriamo l'esecuzione \(D(d)\), cioè
il programma \(D\) applicato alla propria descrizione. Si presentano
due possibilità:

\begin{enumerate}
  \item Se \(O(d,d)=1\), l'osservatore prevede che \(D(d)\) termini.
        Per costruzione, però, \(D\) esegue un ciclo senza fine,
        contraddicendo la previsione.

  \item Se \(O(d,d)=0\), l'osservatore prevede che \(D(d)\) non termini.
        Per costruzione, invece, \(D\) termina immediatamente,
        contraddicendo nuovamente la previsione.
\end{enumerate}

Entrambe le risposte possibili conducono a una contraddizione.
L'ipotesi dell'esistenza di un osservatore algoritmico sempre
terminante e corretto è quindi falsa: il problema dell'arresto
non è decidibile in generale.

La contraddizione nasce dalla possibilità di trattare la descrizione
di un programma come un dato, consentendo di applicare \(D\)
alla propria descrizione. Non è necessario che il programma modifichi
il proprio codice durante l'esecuzione.

Questo risultato non impedisce di stabilire la terminazione di
programmi particolari. Se un programma termina, per esempio, una
simulazione completa permette di verificarlo. Se invece continua
indefinitamente, la sola simulazione non permette di concludere
in un tempo finito che non terminerà mai. La
\cref{prop:fi-problema-arresto} esclude un algoritmo capace di
risolvere correttamente tutti i casi possibili, non la possibilità
di analizzarne singoli casi mediante altri ragionamenti.

Il problema dell'arresto mostra dunque un limite intrinseco del
calcolo algoritmico, indipendente dalla potenza della macchina
utilizzata. Il nastro illimitato della macchina di Turing è infatti
un'idealizzazione: un computer reale dispone di memoria finita,
tempi di esecuzione limitati e risorse fisiche. Può simulare passi
del modello entro le risorse disponibili, ma non coincide letteralmente
con una macchina a memoria illimitata. Un limite di calcolabilità
non va pertanto confuso con un'insufficienza di memoria o con un
tempo di esecuzione eccessivo.

\section{Architettura del calcolatore ed esecuzione delle istruzioni}
\label{sec:fi-architettura-calcolatore}

Il modello della macchina di Turing chiarisce i limiti teorici del calcolo;
per comprendere come i programmi vengano eseguiti dai calcolatori reali
occorre invece esaminarne l'architettura. Questa descrive l'organizzazione
delle componenti hardware e le modalità con cui cooperano.
Tra i modelli classici si distinguono l'\emph{architettura di von Neumann}
e l'\emph{architettura Harvard}.

Nel modello di von Neumann, la CPU, la memoria centrale e i dispositivi
di ingresso e uscita comunicano attraverso collegamenti condivisi.
Istruzioni e dati risiedono nella stessa memoria e utilizzano il
medesimo percorso di accesso. Il \emph{principio del programma memorizzato}
prevede che anche le istruzioni siano conservate in memoria, consentendo
di cambiare il programma eseguito senza modificare i circuiti della
macchina. Tale principio non è esclusivo dell'architettura di von Neumann.

L'architettura Harvard distingue invece le memorie e i percorsi di
accesso destinati alle istruzioni da quelli riservati ai dati,
permettendo di effettuare i due tipi di accesso contemporaneamente.
Le architetture moderne possono combinare caratteristiche di entrambi
i modelli: per esempio, molti processori dispongono di uno spazio di
indirizzamento unificato e di cache separate per istruzioni e dati.

La \emph{CPU} (\emph{Central Processing Unit}) preleva ed esegue le
istruzioni. Al suo interno, l'\emph{unità di controllo} interpreta le
istruzioni e coordina le operazioni richieste, mentre l'\emph{unità
logico-aritmetica}, o \emph{ALU} (\emph{Arithmetic Logic Unit}), esegue
operazioni aritmetiche e logiche, come somme e confronti.
I \emph{registri} sono piccole memorie interne alla CPU, rapidamente
accessibili, che conservano temporaneamente dati, indirizzi e altre
informazioni necessarie all'esecuzione.

La \cref{fig:fi-von-neumann} rappresenta i principali blocchi funzionali
del modello di von Neumann.

\begin{figure}[htbp]
  \centering
  % Diagramma originale; licenza CC BY-SA 4.0 del materiale didattico.
\begin{tikzpicture}[unipd/diagram]
  \draw[unipd/edge,rounded corners=2pt] (-4.6,0.5) rectangle (4.6,2.7);
  \node at (0,2.4) {CPU};
  \node[unipd/architecture component] (control) at (-3,1.35)
    {Unità di\\controllo};
  \node[unipd/architecture component] (alu) at (0,1.35) {ALU};
  \node[unipd/architecture component] (registers) at (3,1.35)
    {Registri\\PC, istruzione\\dati};
  \draw[unipd/edge,<->] (control) -- (alu);
  \draw[unipd/edge,<->] (alu) -- (registers);
  \draw[unipd/edge,line width=1.2pt] (-4.6,-0.3) -- (4.6,-0.3);
  \node[font=\normalfont,anchor=north] at (0,-0.4)
    {Bus: dati, indirizzi, controllo};
  \draw[unipd/edge,<->] (0,0.5) -- (0,-0.3);
  \node[unipd/architecture component] (memory) at (-3,-1.5)
    {Memoria centrale\\istruzioni e dati};
  \node[unipd/architecture component] (io) at (3,-1.5)
    {Interfacce I/O\\periferiche};
  \draw[unipd/edge,<->] (memory.north) -- (-3,-0.3);
  \draw[unipd/edge,<->] (io.north) -- (3,-0.3);
\end{tikzpicture}

  \caption{Organizzazione funzionale di von Neumann. Il collegamento
    condiviso comprende segnali di dati, indirizzi e controllo;
    le frecce indicano lo scambio complessivo, non la direzione di ogni segnale.}
  \label{fig:fi-von-neumann}
\end{figure}

Le comunicazioni fra le componenti sono rappresentate mediante
\emph{bus}, ossia canali che consentono il trasferimento di informazioni
secondo regole condivise. Si distinguono funzionalmente:
\begin{itemize}
  \item il \emph{bus dei dati}, che trasferisce dati e istruzioni;
  \item il \emph{bus degli indirizzi}, che identifica la locazione di
        memoria o il dispositivo interessato;
  \item il \emph{bus di controllo}, che trasporta i segnali necessari
        a coordinare operazioni come lettura e scrittura.
\end{itemize}
La distinzione riguarda la funzione dei segnali e non implica
necessariamente la presenza di tre collegamenti fisici separati.
Per leggere un valore dalla memoria, per esempio, si specifica
l'indirizzo, si richiede l'operazione mediante i segnali di controllo
e si trasferisce il risultato attraverso il canale dei dati.

La condivisione dei collegamenti semplifica l'organizzazione del sistema,
ma introduce un limite alle prestazioni: più richieste possono competere
per la stessa risorsa. La \emph{banda}, cioè la quantità di dati
trasferibile per unità di tempo, è limitata, e la CPU può dover attendere
il completamento degli accessi alla memoria. Questo fenomeno è alla
base del cosiddetto \emph{collo di bottiglia di von Neumann}.
Il modello rappresentato è una descrizione funzionale; i calcolatori
moderni possono utilizzare interconnessioni più articolate.

Il registro \emph{program counter} (PC) contiene l'indirizzo della
prossima istruzione da prelevare. Nel modello didattico, il
\emph{registro istruzione} conserva l'istruzione corrente durante
la sua interpretazione. L'esecuzione segue un ciclo fondamentale,
articolato nelle seguenti fasi:
\begin{enumerate}
  \item \emph{fetch}: l'istruzione individuata dal PC viene prelevata
        dalla memoria;
  \item \emph{decode}: l'unità di controllo identifica l'operazione
        e gli operandi richiesti;
  \item \emph{execute}: vengono svolte le operazioni previste
        dall'istruzione, utilizzando gli operandi necessari;
  \item \emph{write-back}: quando previsto, il risultato viene scritto
        nella destinazione indicata, per esempio un registro o la memoria.
\end{enumerate}

Durante l'esecuzione sequenziale, il PC viene aggiornato all'indirizzo
dell'istruzione successiva, tenendo conto della sua lunghezza.
Un'istruzione di salto può invece modificarlo, eventualmente in base
all'esito di una condizione, permettendo di realizzare selezioni e cicli.
Non tutte le istruzioni richiedono una scrittura del risultato.

Il \emph{clock} è un segnale periodico che scandisce il funzionamento
delle componenti sincrone. La sua frequenza, misurata in hertz,
indica il numero di cicli al secondo. Un ciclo di clock non coincide
necessariamente con l'esecuzione completa di un'istruzione: la frequenza
da sola non determina quindi le prestazioni di un processore.

Per migliorare le prestazioni si può sfruttare il
\emph{parallelismo a livello di istruzioni}
(ILP, \emph{Instruction-Level Parallelism}), sovrapponendo
l'elaborazione di istruzioni diverse.

Nella \emph{pipeline}, l'esecuzione viene suddivisa in stadi.
Mentre un'istruzione attraversa uno stadio, altre istruzioni possono
occupare gli stadi precedenti o successivi. Una suddivisione didattica
a cinque stadi comprende prelievo, decodifica, recupero degli operandi,
esecuzione e scrittura del risultato, come illustrato nella
\cref{fig:fi-pipeline}.

\begin{figure}[htbp]
  \centering
  % Diagramma originale; licenza CC BY-SA 4.0 del materiale didattico.
\begin{tikzpicture}[unipd/diagram,x=0.9cm,y=0.65cm]
  \node[anchor=east] at (-0.7,0) {Ciclo};
  \foreach \cycle in {1,...,8} {
    \node at (\cycle,0) {\cycle};
  }
  \foreach \instruction in {1,...,4} {
    \node[anchor=east] at (-0.7,-\instruction) {Istruzione \instruction};
    \foreach \stage/\name in {0/F,1/D,2/O,3/E,4/W} {
      \node[unipd/node,rounded corners=0pt,minimum width=8mm,
        minimum height=6mm,inner sep=0pt]
        at ({\instruction+\stage},-\instruction) {\name};
    }
  }
\end{tikzpicture}

  \caption{Pipeline ideale a cinque stadi, ciascuno della durata di un ciclo,
    senza attese: F = prelievo, D = decodifica, O = recupero operandi,
    E = esecuzione, W = scrittura. Dal quinto ciclo viene completata
    un'istruzione per ciclo finché la pipeline rimane alimentata.}
  \label{fig:fi-pipeline}
\end{figure}

La pipeline può aumentare il numero di istruzioni completate per unità
di tempo, senza necessariamente ridurre la durata di una singola
istruzione. Dipendenze fra istruzioni e attese per le risorse possono
inoltre impedire il funzionamento al ritmo ideale.

Un processore \emph{superscalare} dispone di più unità funzionali e può
avviare più istruzioni nello stesso ciclo, purché le dipendenze e
le risorse disponibili lo consentano. La pipeline sovrappone stadi
di istruzioni diverse, mentre l'organizzazione superscalare consente
di sfruttare più unità funzionali nello stesso intervallo di tempo.
Le due tecniche possono coesistere.

Un'ulteriore forma di parallelismo riguarda l'impiego di più
\emph{core} o processori. Un processore \emph{multicore} integra
più core capaci di eseguire flussi di istruzioni; un sistema
multiprocessore comprende più processori, eventualmente con memoria
condivisa. Nei \emph{multicalcolatori}, invece, nodi dotati di memoria
propria comunicano mediante messaggi. In tutti questi casi,
l'accelerazione dipende dalla possibilità di distribuire il lavoro
e dai costi di coordinamento: aggiungere unità di calcolo non rende
automaticamente più veloce un procedimento sequenziale.

La \emph{GPU} (\emph{Graphics Processing Unit}), originariamente
sviluppata per la grafica, è adatta anche all'elaborazione parallela
di grandi quantità di dati. Rispetto a una CPU, progettata per gestire
in modo efficiente flussi di controllo complessi e operazioni con
bassa latenza, una GPU privilegia l'esecuzione di numerose operazioni
simili su dati differenti.

Elaborazione di immagini, operazioni su matrici e simulazioni sono
esempi di applicazioni che possono beneficiare di questa organizzazione.
Il vantaggio dipende tuttavia dalla quantità di lavoro parallelizzabile,
dalle dipendenze fra operazioni e dai trasferimenti dei dati.
Le unità di calcolo di una GPU non sono direttamente equivalenti
ai core di una CPU.

\emph{CUDA} (\emph{Compute Unified Device Architecture}) è una
piattaforma NVIDIA per la programmazione del calcolo su GPU.
In questo modello la CPU, detta \emph{host}, coordina il programma
e avvia sulla GPU, detta \emph{device}, funzioni di calcolo chiamate
\emph{kernel}.

L'esecuzione di un kernel coinvolge numerosi \emph{thread}, organizzati
in \emph{blocchi}; l'insieme dei blocchi costituisce una \emph{griglia}.
Il modello SIMT (\emph{Single Instruction, Multiple Threads}) permette
a gruppi di thread di eseguire istruzioni su dati differenti.
Questo non implica che tutti i thread vengano eseguiti fisicamente
nello stesso istante. CPU e GPU cooperano nell'elaborazione:
CUDA è una piattaforma di programmazione, non un algoritmo né l'unica
tecnologia disponibile per il calcolo su GPU.\footnote{NVIDIA,
\href{https://docs.nvidia.com/cuda/cuda-programming-guide/01-introduction/programming-model.html}
{\emph{CUDA Programming Guide}, § 1.2, Programming Model},
in particolare sistemi eterogenei, blocchi, griglie e SIMT.}

\section{Memoria, archiviazione e dispositivi di ingresso e uscita}
\label{sec:fi-memoria-archiviazione}

L'esecuzione delle istruzioni dipende anche dal modo in cui dati e
programmi vengono memorizzati e trasferiti. Nel contesto della
memorizzazione, un \emph{bit} indica una cifra binaria, con valori
convenzionali \(0\) e \(1\); un \emph{byte} è un gruppo di otto bit.
Il bit come unità di informazione, introdotto nella
\cref{sec:fi-informatica-calcolo}, è un concetto distinto.
Queste unità descrivono la quantità di dati memorizzati, non il loro
significato, che dipende dalle convenzioni illustrate nel
\cref{chap:rappresentazione-informazioni}.

La memoria è organizzata in \emph{locazioni}, ciascuna identificata
da un \emph{indirizzo}. Un'operazione di lettura ne restituisce il
contenuto, mentre una scrittura lo sostituisce con il dato fornito.
Nelle architetture indirizzabili per byte, ogni indirizzo identifica
un byte, anche se una singola operazione può trasferirne più di uno.
L'unità minima indirizzabile non coincide quindi necessariamente
con la quantità di dati trasferiti.

Lo \emph{spazio di indirizzamento} è l'insieme degli indirizzi disponibili.
Se gli indirizzi sono rappresentati mediante \(a\) bit, si possono
distinguere al più \(2^a\) locazioni. Nel caso di memoria indirizzabile
per byte, ciò corrisponde a \(2^a\) byte: con sedici bit si ottengono
\(2^{16}\) byte, ossia \(64\) KiB, mentre con trentadue bit si raggiungono
\(2^{32}\) byte, equivalenti a \(4\) GiB.

La capacità teoricamente indirizzabile non coincide necessariamente
con la RAM installata: alcune regioni dello spazio di indirizzamento
possono essere riservate ad altri dispositivi o non essere utilizzabili.
Inoltre, la dimensione dei registri di un processore non determina
necessariamente il numero di bit effettivamente impiegati per gli
indirizzi di memoria.

Per esprimere le capacità occorre distinguere i prefissi decimali
da quelli binari. Il simbolo B indica il byte:
\(1\) kB corrisponde a \(10^3\) B, \(1\) MB a \(10^6\) B e
\(1\) GB a \(10^9\) B. I prefissi binari utilizzano invece potenze
di due: \(1\) KiB vale \(2^{10}\) B, \(1\) MiB vale \(2^{20}\) B
e \(1\) GiB vale \(2^{30}\) B. Le due convenzioni non sono
intercambiabili e devono essere distinte quando si indicano
capacità di memoria.\footnote{NIST,
\href{https://physics.nist.gov/cuu/Units/binary.html}
{\emph{Prefixes for binary multiples}},
prefissi IEC e confronto con i prefissi SI.}

La \emph{memoria centrale} contiene dati e istruzioni necessari
all'esecuzione dei programmi. La memoria \emph{RAM}
(\emph{Random Access Memory}), utilizzata come memoria di lavoro,
consente operazioni di lettura e scrittura ed è normalmente
\emph{volatile}: il suo contenuto viene perso quando manca
l'alimentazione. L'espressione «accesso casuale» indica che una
locazione può essere raggiunta direttamente, senza dover attraversare
quelle precedenti, e non che gli indirizzi vengano scelti casualmente.
Ciò non implica che tutti gli accessi abbiano necessariamente
la stessa latenza.

La memoria \emph{ROM} (\emph{Read-Only Memory}) è invece non volatile
ed è destinata principalmente alla lettura durante il normale
funzionamento. Tradizionalmente viene utilizzata per conservare
il \emph{firmware}, ossia il software strettamente legato al
funzionamento di un dispositivo, comprese le procedure di avvio.
Nei sistemi moderni il firmware può risiedere anche in memorie
non volatili riprogrammabili, come quelle flash.
La possibilità di modificare il contenuto e la capacità di
conservarlo senza alimentazione sono quindi proprietà distinte.

La \emph{memoria secondaria}, o \emph{memoria di massa}, permette
di conservare dati e programmi anche quando non sono in esecuzione.
Rispetto alla memoria centrale offre generalmente capacità maggiori,
ma tempi di accesso più elevati. Tra i dispositivi più comuni si
distinguono i dischi rigidi e le unità a stato solido.

Un \emph{disco rigido} (HDD, \emph{Hard Disk Drive}) memorizza
i dati magneticamente su piatti rotanti; l'accesso richiede il
posizionamento delle testine e può comportare attese dovute alla
rotazione. Un'\emph{unità a stato solido} (SSD, \emph{Solid State Drive})
utilizza normalmente memoria flash, priva di parti meccaniche mobili,
e offre in genere tempi di accesso inferiori. Entrambi i dispositivi
sono non volatili, ma rimangono soggetti a guasti e perdita dei dati.

Durante l'esecuzione, i dati necessari vengono trasferiti dalla
memoria secondaria alla memoria centrale. Le modifiche presenti
nella memoria di lavoro devono invece essere salvate su un
supporto persistente per poter essere conservate dopo lo spegnimento.

Le differenze di velocità, capacità e costo motivano
l'organizzazione in una \emph{gerarchia di memoria}, che comprende
registri, cache, memoria centrale e memoria secondaria.
I livelli più vicini alla CPU offrono generalmente accessi più
rapidi, ma capacità inferiori rispetto a quelli destinati
all'archiviazione.

La \emph{cache} è una memoria piccola e veloce che conserva copie
di dati e istruzioni per ridurre la necessità di accedere a memorie
più lente. La sua efficacia dipende dal comportamento dei programmi,
che spesso presenta due forme di \emph{località}:
\begin{itemize}
  \item \emph{località temporale}: un dato utilizzato recentemente
        ha una maggiore probabilità di essere utilizzato nuovamente;
  \item \emph{località spaziale}: l'accesso a una locazione è spesso
        seguito dall'accesso a locazioni vicine.
\end{itemize}
Queste proprietà permettono alla cache di migliorare le prestazioni
quando i dati richiesti sono già disponibili al suo interno.
Non costituiscono tuttavia garanzie valide per ogni programma
o sequenza di accessi.

I dispositivi di \emph{ingresso e uscita} (I/O, \emph{Input/Output})
consentono lo scambio di dati fra il calcolatore e l'ambiente esterno.
Una tastiera acquisisce informazioni, mentre uno schermo presenta
i risultati dell'elaborazione. Altri dispositivi, come schede di rete
e unità di archiviazione, svolgono entrambe le funzioni.
Le periferiche possono comunicare anche con altri sistemi,
senza richiedere necessariamente l'interazione di un utente.

Il trasferimento dei dati viene coordinato attraverso interfacce
e meccanismi di controllo differenti.

Nell'\emph{I/O programmato} la CPU esegue direttamente le istruzioni
necessarie al trasferimento. Può inoltre verificare ripetutamente
lo stato della periferica mediante \emph{polling}, occupando
tempo di elaborazione durante l'attesa.

Gli \emph{interrupt} permettono invece a un dispositivo di segnalare
un evento che richiede attenzione. La CPU può interrompere il
flusso corrente, conservare le informazioni necessarie a riprenderlo
ed eseguire una procedura di servizio. Terminata la gestione
dell'evento, può riprendere il lavoro precedentemente sospeso.

Il \emph{DMA} (\emph{Direct Memory Access}) consente a un controller
di trasferire blocchi di dati fra memoria e periferiche senza
richiedere che la CPU esegua ogni singola operazione di trasferimento.
La CPU configura l'operazione e può ricevere un interrupt quando
questa è completata.

Polling, interrupt e DMA svolgono dunque funzioni differenti e
possono essere utilizzati in combinazione. In particolare,
il DMA riduce il coinvolgimento diretto della CPU nei trasferimenti,
ma richiede comunque il coordinamento nell'accesso alla memoria.

\section{Linguaggi di programmazione e forme di esecuzione}
\label{sec:fi-linguaggi-esecuzione}

L'organizzazione dell'hardware chiarisce come vengono eseguite le
operazioni, ma non come queste siano specificate in un programma.
I linguaggi naturali permettono di esprimere concetti il cui significato
può dipendere dal contesto e risultare ambiguo. Per descrivere istruzioni
che un calcolatore possa interpretare secondo regole precise, si
utilizzano invece linguaggi formali.

Un \emph{linguaggio formale} è un insieme di stringhe costruite a
partire da un \emph{alfabeto}, cioè un insieme finito di simboli.
Una \emph{stringa} è una sequenza finita di simboli; la stringa vuota
si indica con \(\varepsilon\). Se \(\Sigma\) è un alfabeto,
\(\Sigma^*\) rappresenta l'insieme di tutte le stringhe finite
costruibili sui suoi simboli, inclusa quella vuota. Un linguaggio
formale è pertanto un sottoinsieme \(L\subseteq\Sigma^*\).

Una \emph{grammatica formale} descrive come generare le stringhe
di un linguaggio mediante simboli terminali, simboli non terminali,
un simbolo iniziale e regole di produzione. I terminali costituiscono
i simboli delle stringhe prodotte, mentre i non terminali vengono
utilizzati durante la derivazione.

Per esempio, considerando i terminali \(a,b\), il simbolo iniziale
non terminale \(S\) e le produzioni
\[
  S\to aSb,
  \qquad S\to ab,
\]
si ottengono le stringhe \(ab,aabb,aaabbb,\ldots\).
La prima regola aggiunge un simbolo \(a\) a sinistra e un simbolo
\(b\) a destra; la seconda conclude la derivazione. Il linguaggio
generato contiene quindi stringhe con lo stesso numero positivo di
\(a\) e \(b\), disposti in quest'ordine.

Noam Chomsky ha classificato le grammatiche in una gerarchia basata
sulle restrizioni imposte alle produzioni. In ordine di potere
espressivo crescente si distinguono:
\begin{itemize}
  \item i \emph{linguaggi regolari}, riconosciuti da automi finiti;
  \item i \emph{linguaggi liberi dal contesto}, riconosciuti da
        automi a pila;
  \item i \emph{linguaggi sensibili al contesto}, riconosciuti da
        macchine di Turing con memoria limitata linearmente
        rispetto alla dimensione dell'ingresso;
  \item i \emph{linguaggi di tipo zero}, riconosciuti da macchine
        di Turing generali.
\end{itemize}
Ogni classe è contenuta nella successiva. La gerarchia mette in
relazione la struttura delle grammatiche con le capacità dei modelli
di calcolo. Occorre tuttavia distinguere il \emph{riconoscimento}
dalla \emph{decisione}: una macchina di Turing che riconosce un
linguaggio di tipo zero può non terminare quando riceve una stringa
che non vi appartiene, mentre un procedimento decisore deve
terminare per ogni ingresso.

Le grammatiche permettono di descrivere la struttura sintattica
dei linguaggi di programmazione, ma non ne determinano da sole
il significato e il comportamento.\footnote{Per la gerarchia
di Chomsky: J. E. Hopcroft, R. Motwani e J. D. Ullman,
\emph{Introduction to Automata Theory, Languages, and Computation},
terza edizione, Pearson.}

Un processore esegue direttamente operazioni appartenenti al proprio
\emph{insieme di istruzioni} (\emph{instruction set}), che ne definisce
le operazioni disponibili e le relative codifiche. Il
\emph{linguaggio macchina} esprime i programmi attraverso tali
codifiche binarie. Le istruzioni possono specificare, per esempio,
trasferimenti di dati, operazioni aritmetiche, confronti e salti.
Poiché gli insiemi di istruzioni dipendono dall'architettura,
un programma in linguaggio macchina non è necessariamente
eseguibile su processori differenti.

Il \emph{linguaggio assembly} rappresenta le istruzioni macchina
mediante nomi mnemonici e permette di utilizzare riferimenti
simbolici per dati e indirizzi. Un programma detto
\emph{assemblatore} traduce il codice assembly in codice macchina.
Il linguaggio rimane strettamente legato all'architettura del
processore, anche se non ogni riga del sorgente deve necessariamente
corrispondere a una singola istruzione macchina: possono essere
presenti direttive e altre costruzioni elaborate dall'assemblatore.

I \emph{linguaggi ad alto livello}, come Python e Java, permettono
invece di descrivere elaborazioni e strutture di controllo senza
specificare direttamente le operazioni sui registri del processore.
Una singola istruzione del linguaggio può richiedere numerose
istruzioni macchina, rendendo necessaria una forma di traduzione
o un ambiente capace di eseguirla.

Il testo di un programma nel linguaggio scelto costituisce il
\emph{codice sorgente}. La \emph{sintassi} stabilisce quali costruzioni
sono ammesse, mentre la \emph{semantica} ne definisce il significato
e il comportamento. Un programma può rispettare tutte le regole
sintattiche senza essere corretto rispetto al problema che deve
risolvere.

Per trasformare il codice sorgente in operazioni eseguibili si
utilizzano differenti tecniche. Un \emph{compilatore} traduce un
programma in un'altra rappresentazione, che può essere direttamente
codice macchina oppure \emph{codice intermedio} destinato a
ulteriori elaborazioni.

Quando la compilazione produce codice macchina destinato
all'esecuzione nativa, può essere necessaria una fase di
\emph{collegamento} (\emph{linking}) con altri moduli o librerie.
Nei sistemi dotati di sistema operativo, l'eseguibile viene quindi
caricato e avviato. Deve essere compatibile sia con l'architettura
del processore sia con l'ambiente di esecuzione:
la portabilità del codice sorgente non garantisce automaticamente
quella del programma compilato.

Un \emph{interprete} è invece un programma che realizza le
operazioni previste da un linguaggio, elaborando direttamente
il sorgente oppure una sua rappresentazione intermedia.
L'interprete viene a sua volta eseguito dal processore e non
ne sostituisce il funzionamento.

\emph{Compilazione} e \emph{interpretazione} non sono necessariamente
alternative: descrivono modalità di elaborazione ed esecuzione che
possono essere combinate all'interno dello stesso ambiente.

Per esempio, nell'implementazione di Python denominata
\emph{CPython}, il codice sorgente viene tradotto in una
rappresentazione intermedia chiamata \emph{bytecode}, successivamente
eseguita dall'interprete. Nel percorso usuale di Java, il
compilatore produce bytecode destinato alla
\emph{Java Virtual Machine} (JVM).

La JVM è un ambiente di esecuzione software, non un processore
fisico: può interpretare il bytecode e tradurne alcune parti in
codice macchina durante l'esecuzione. In entrambi i casi,
fra il sorgente e le istruzioni effettivamente eseguite dalla CPU
possono quindi esistere più livelli di rappresentazione e
traduzione.\footnote{Per la distinzione fra rappresentazioni ed
esecuzione: documentazione Python,
\href{https://docs.python.org/3/glossary.html\#term-bytecode}
{voce \emph{bytecode}}; Oracle,
\emph{The Java Virtual Machine Specification},
\href{https://docs.oracle.com/javase/specs/jvms/se25/html/jvms-1.html\#jvms-1.2}
{§ 1.2, The Java Virtual Machine}.}